\subsubsection{Constructive determination and role of the dodecaphasic register}
\label{alpha-ampl-dependencias}

The preceding definition clarifies what it means to determine alpha from
the structure. The domain of the operation contains four registers of
known provenance: the three regional sequences and the integer phase
register. Normalization also incorporates a boundary condition
compatible with extension. Once these data are fixed, Euclidean
division determines each remainder and its quotient. Proposition
\ref{prop:alpha-normalizacion} extends that determination to the complete
system. Independence from a target physical value is therefore
a property of the order of construction: the value compared at the
end does not participate in selecting any of the arguments of this
map.

This causal independence must be distinguished from algebraic independence
between numbers. The former is established by identifying the domain, the
operations and the provenance of the coefficients; the latter would be a
statement about polynomial relations. The result used here is the
former. The construction determines a coordinate through registers
produced by APP--TRIT--TPK and allows the relations it satisfies to be
checked subsequently. This order makes it possible to use a coordinate already
generated in a later operation without turning it into an external parameter.

The register $K$ has a more precise role than that of a numerical
correction added to a formula. Its twelve positions arise from the
reversible transformation of the signed register. The equation
$U=\mathcal H_{12}K$ preserves the possibility of returning to that register,
while the transverse differences and the total sum allow another
reconstruction. In the periodic evaluation, the order of the positions
determines the positional weight of each component. The same set of
components, arranged in a different order, generally produces a different rational
coordinate. Thus, the phase origin and orientation form part of the
definition of the number entering \eqref{eq:alpha-identidad-completa}.

The normalization of alpha composes this information with the regions. Its
local identity distinguishes two successive operations. First the oriented
balance of blocks is formed; then that balance is represented through
a canonical remainder and a quotient linking adjacent positions. The first
operation preserves the sign of the balance and allows a support to be
extracted subsequently. The second allows the prefixes of the real
coordinate to be formed. Both readings remain related because the state preserves the
data before and after normalization. The signed support and the
positional expansion are thus different realizations of an explicit
arithmetic composition.

The telescoping identity expresses the global content of this composition.
The interior quotients cancel in pairs when a prefix is evaluated,
while their endpoints remain in the equality. This cancellation is a
property of the weighted sum; it neither removes the quotients from the state nor
allows them to be suppressed before the sum is taken. The right boundary of a
window receives the contribution of the continuation. For this reason,
evaluating an isolated window and restricting an extended
section are distinct operations. The change of the final block from $663$ to
$662$, proved above, provides an explicit realization of this
difference without altering the blocks already stabilized to its left.

In this sense, bidirectionality has a concrete arithmetic
formulation. Generation provides compatible extensions;
normalization transmits the integer information of their boundary from
positions of lower weight towards those of higher weight. The equation preserves the
contribution of both directions through its indices and its boundary
conditions. Reversibility in the fixed fiber is checked by reconstructing
$K$ from the other registers and the complete quotients. This
inversion certifies the consistency of the original operation; the genealogical
direction remains the one that produces the register first and the alpha
coordinate afterwards.

The second determination has complementary content. In it, the
immediate object is a balance of vacancies, and the parameter appears as
its root in an isolating domain. The positional route provides a
canonical normalization of registers; the analytic route studies the vanishing
of an operator. Presenting both requires precisely this
difference of functions to be preserved. Their relations with the common state must be
expressed through the refinement maps and the comparison of
their complete realizations, as developed in the corresponding
subsection. Agreement of a finite approximation constitutes a
localized check; identification of the complete constructions
has the domain and scope of its own statement.

Finally, this organization explains the position of alpha in the article.
Its equation appears after regional generation and the construction
of $K$, because it uses both results. The electron section then uses
that coordinate as a factor in its composition. The exceptional
extension receives the registers and their supports, where incidence
information is preserved. These two continuations share the same
provenance and develop distinct mathematical functions: evaluation of
a central section and construction of a boundary configuration.
Developing each one allows the structure to be followed through the
equation, instead of restricting the comparison to its scalar value.
